% ============================================================
% CHƯƠNG 1: GIỚI THIỆU
% MỤC ĐÍCH: Đáp ứng Câu 1 (mục tiêu đề tài) + Câu 2 (yêu cầu đề tài)
% PHÂN BỔ: ~8–10 trang
% ============================================================
\chapter{Giới thiệu}
\label{chap:gioi-thieu}

\section{Bối cảnh và động lực nghiên cứu}
\label{sec:boi-canh}

Trong nửa sau của thập kỷ 2020, sự bùng nổ của các mô hình ngôn ngữ lớn
(\textit{Large Language Models} -- LLM) như GPT-4, LLaMA-3, Gemini và Claude đã
đẩy nhu cầu tính toán trên bộ tăng tốc đồ họa (GPU) tăng theo cấp số nhân. Các
phiên huấn luyện một mô hình tiên tiến hiện nay đòi hỏi hàng nghìn GPU NVIDIA
H100 hoặc A100 chạy liên tục trong nhiều tuần, kéo theo sự mở rộng quy mô chưa
từng có của các trung tâm dữ liệu (\textit{datacenter}) chuyên dụng cho trí tuệ
nhân tạo (AI)~\cite{patterson2021carbon, luccioni2023estimating}. Theo báo cáo
\textit{Electricity 2024} của Cơ quan Năng lượng Quốc tế 
(IEA)~\cite{iea2024electricity}, tổng năng lượng tiêu thụ của các trung tâm dữ
liệu, AI và tiền mã hóa trên toàn cầu trong năm 2022 ước đạt khoảng
$460$~TWh -- tương đương $1{,}5\%$ tổng nhu cầu điện năng thế giới -- và dự báo
sẽ tăng gấp đôi, chạm ngưỡng $1000$~TWh vào năm 2026, vượt qua tổng tiêu thụ
điện của cả nước Nhật Bản.

\begin{figure}[H]
  \centering
  \includegraphics[width=0.85\textwidth]{reference/dc-energy-growth.png}
  % \fbox{\parbox[c][4.5cm][c]{0.8\textwidth}{\centering
  %   \textit{[Hình 1.1 -- Tăng trưởng năng lượng tiêu thụ datacenter toàn cầu
  %   2018--2026, nguồn IEA Electricity 2024]}}}
  \caption{Xu hướng tiêu thụ điện năng của trung tâm dữ liệu, AI và tiền mã hóa
  trên toàn cầu giai đoạn 2018--2026. Nguồn: IEA~\cite{iea2024electricity}.}
  \label{fig:dc-energy-growth}
\end{figure}

Tuy nhiên, cùng với tốc độ mở rộng đó là một nghịch lý đáng lo ngại về hiệu
suất sử dụng tài nguyên. Phân tích vết tải (\textit{trace}) công khai của
Alibaba PAI Cluster v2023~\cite{weng2023alibaba} cho thấy: mặc dù chi phí GPU
chiếm tới hơn $80\%$ tổng giá thành phần cứng của một AI cluster, hiệu suất sử
dụng GPU trung bình chỉ đạt $\approx 25$--$40\%$. Nguyên nhân chính nằm ở hiện
tượng \textbf{phân mảnh tài nguyên} (\textit{Resource Fragmentation}): một host
có thể còn dư GPU nhưng không còn đủ CPU/RAM để nhận thêm tác vụ, hoặc ngược
lại, dẫn đến những GPU ``mồ côi'' (\textit{orphan GPUs}) tiêu tốn năng lượng ở
trạng thái \textit{idle} mà không sinh giá trị tính
toán~\cite{weng2022mlaas, jeon2019analysis}.


Bài toán càng trở nên phức tạp khi nhà vận hành phải đồng thời theo đuổi
\textbf{hai mục tiêu xung đột}:
\begin{enumerate}
  \item \textbf{Tiết kiệm năng lượng} -- giảm tổng điện năng tiêu thụ và lượng
        khí thải $\mathrm{CO_2}$ tương ứng;
  \item \textbf{Đảm bảo Thỏa thuận Mức Dịch vụ} (\textit{Service Level
        Agreement} -- SLA) -- hoàn thành tác vụ trước hạn chót, đặc biệt là các
        tác vụ \textit{LatencySensitive}.
\end{enumerate}
Hai mục tiêu này thường đối lập: việc \textit{consolidation} tác vụ vào ít host
hơn để tắt bớt các host nhàn rỗi giúp tiết kiệm điện, nhưng lại làm tăng độ trễ
xếp hàng và nguy cơ vi phạm SLA. Ngược lại, phân tán tải để đảm bảo SLA sẽ
khiến nhiều host vận hành ở mức tải thấp.

Trong bối cảnh Việt Nam, khi giá điện công nghiệp liên tục được điều chỉnh tăng
và các cam kết Net-Zero 2050 đặt áp lực mạnh lên ngành công nghiệp tính toán,
việc xây dựng một nền tảng nghiên cứu cho phép đánh giá và phát triển các thuật
toán lập lịch \textbf{vừa tiết kiệm năng lượng vừa đảm bảo chất lượng dịch vụ}
trở thành một nhu cầu mang tính cấp thiết và có giá trị thực tiễn rõ rệt.

\section{Phát biểu bài toán}
\label{sec:bai-toan}

\subsection{Đầu vào và đầu ra}
Mô tả bài toán bằng ngôn ngữ kỹ thuật, không hình thức hoá quá mức:

\textbf{Đầu vào.} Hệ thống nhận:
\begin{itemize}
  \item Một cụm gồm $m$ \textbf{host} vật lý. Mỗi host có dung lượng cố định
        về CPU (đơn vị MIPS), RAM (MB), GPU (số chiếc) và băng thông (Mbps).
  \item Một dòng \textbf{tác vụ} đến theo thời gian. Mỗi tác vụ $\tau_i$
        khai báo nhu cầu tài nguyên $(r^{\mathrm{cpu}}_i,
        r^{\mathrm{ram}}_i, r^{\mathrm{gpu}}_i)$, thời hạn hoàn thành $d_i$
        và nhãn chất lượng dịch vụ $q_i \in \{\text{Burstable},
        \text{LatencySensitive}, \text{BestEffort}\}$ -- lấy từ Alibaba
        trace.
\end{itemize}

\textbf{Đầu ra.} Một quyết định: \emph{tác vụ $\tau_i$ được đặt lên host
nào, hay tạm hoãn lại đến bước kế tiếp.} Ký hiệu phép gán này là
$f(\tau_i)$.

\textbf{Ràng buộc bắt buộc.} Tại mọi thời điểm, tổng nhu cầu tài nguyên
của các tác vụ đang chạy trên một host không được vượt quá dung lượng
của host đó:
\begin{equation}
  \text{Tổng } r^{(k)}_i \text{ trên host } h_j \;\le\; c^{(k)}_j,
  \quad k \in \{\mathrm{cpu},\mathrm{ram},\mathrm{gpu}\}.
  \label{eq:capacity}
\end{equation}
Đây là điều kiện cứng -- vi phạm là không khả thi, không phải chỉ ``tệ
hơn''. Hệ thống thực hiện ràng buộc này qua cơ chế \textbf{action
masking} (Chương~\ref{chap:kien-truc}) -- không cho phép RL agent chọn
host không đủ tài nguyên ngay từ không gian hành động.

\subsection{Hai mục tiêu cần tối ưu cùng lúc}
\textbf{Mục tiêu 1 -- Tiết kiệm năng lượng.} Mỗi host $h_j$ tại thời
điểm $t$ tiêu thụ công suất $P_j(t)$ phụ thuộc mức sử dụng
$U_j(t) \in [0,1]$. ELDAS dùng \emph{mô hình tuyến tính
SPECpower}~\cite{beloglazov2012energy}: công suất bằng $P^{\mathrm{idle}}$
khi host nhàn rỗi, tăng tuyến tính tới $P^{\mathrm{max}}$ khi
$U_j = 1$:
\begin{equation}
  P_j(U_j) \;=\; P^{\mathrm{idle}}_j
              + U_j \cdot \big(P^{\mathrm{max}}_j - P^{\mathrm{idle}}_j\big).
  \label{eq:power-model}
\end{equation}
\textbf{Tổng năng lượng} của một episode mô phỏng là tổng năng lượng
trên toàn bộ host theo thời gian:
\begin{equation}
  E \;=\; \sum_{\text{host}} \sum_{\text{tick}} P_j(t)\,\Delta t.
  \label{eq:total-energy}
\end{equation}
Đây chỉ là phép \emph{cộng dồn} công suất nhân với độ dài tick theo
thời gian -- không phải phép tích phân giải tích phức tạp. Mục tiêu là
\emph{giảm $E$ càng nhiều càng tốt}.

\textbf{Mục tiêu 2 -- Hoàn thành tác vụ trước hạn (SLA).} Một tác vụ
``trễ hạn'' khi thời điểm hoàn thành $C_i$ vượt deadline $d_i$. Mức độ
trễ được tính là $\max(0,\, C_i - d_i)$ -- bằng $0$ nếu xong sớm hơn
hoặc đúng hạn. Tổng vi phạm trên toàn episode là tổng các mức trễ, có
trọng số theo lớp QoS:
\begin{equation}
  V \;=\; \sum_{\text{task}} \lambda(q_i)\,\max\!\big(0,\, C_i - d_i\big),
  \label{eq:sla-violation}
\end{equation}
\begin{table}[H]
  \centering
  \footnotesize
  \begin{tabular}{lc}
    \toprule
    \textbf{Nhãn QoS} & \textbf{Hệ số phạt $\lambda$} \\
    \midrule
    LatencySensitive (tác vụ tương tác) & $3{,}0$ \\
    Burstable (tác vụ thường) & $1{,}0$ \\
    BestEffort (tác vụ nền) & $0{,}3$ \\
    \bottomrule
  \end{tabular}
  \caption{Hệ số phạt $\lambda$ cho từng lớp QoS -- vi phạm SLA của tác
  vụ tương tác bị phạt nặng gấp 3 lần tác vụ thường.}
  \label{tab:lambda-qos}
\end{table}

Mục tiêu là \emph{giảm $V$ càng nhiều càng tốt}.

\subsection{Vì sao hai mục tiêu này xung đột?}
\label{ssec:tradeoff}

Hai mục tiêu trên không thể đồng thời tối ưu vì chúng kéo scheduler
theo \emph{hai hướng ngược nhau}:

\begin{itemize}
  \item \textbf{Để tiết kiệm năng lượng:} dồn tải vào ít host hơn, tắt
        bớt host nhàn rỗi $\Rightarrow$ host ``bận'' có nguy cơ xếp
        hàng dài, làm tăng trễ $\Rightarrow$ \emph{tăng vi phạm SLA}.
  \item \textbf{Để đảm bảo SLA:} dải tải ra nhiều host hơn, mỗi host
        chỉ chạy vài tác vụ $\Rightarrow$ nhiều host vận hành ở mức tải
        thấp, tiêu hao $P^{\mathrm{idle}}$ trên cả những host gần như
        không làm việc $\Rightarrow$ \emph{tăng năng lượng}.
\end{itemize}

Vì vậy không tồn tại một ``phương án tốt nhất tuyệt đối'' mà có \emph{cả
một dải các phương án đánh đổi}: phương án A tiết kiệm điện hơn nhưng
trễ nhiều hơn, phương án B đảm bảo SLA tốt hơn nhưng tốn điện hơn.

\textbf{Cách so sánh hai phương án.} Quy ước đơn giản: phương án A
\emph{tốt hơn} phương án B nếu A \emph{không tệ hơn} ở cả hai mục tiêu
và \emph{tốt hơn rõ rệt ở ít nhất một}. Nếu A tốt hơn ở mục tiêu này
nhưng tệ hơn ở mục tiêu kia, hai phương án không thể so sánh trực
tiếp -- cả hai đều là \emph{ứng cử viên hợp lệ} mà nhà vận hành có thể
chọn tuỳ chính sách. Tập tất cả các phương án không bị một phương án
nào khác ``tốt hơn'' chính là \textbf{Pareto front} -- đường biên mà
ELDAS cố gắng xấp xỉ.

\begin{figure}[H]
  \centering
  \fbox{\parbox[c][4.5cm][c]{0.7\textwidth}{\centering
    \textit{[Hình 1.3 -- Đồ thị 2D minh họa trade-off Energy vs. SLA và đường
    Pareto front]}}}
  \caption{Trade-off giữa năng lượng và vi phạm SLA. Mỗi điểm là một chính sách
  lập lịch; đường cong nối các điểm không bị thống trị (Pareto front) là kết
  quả mà Phase 2 hướng tới xấp xỉ.}
  \label{fig:pareto-illustration}
\end{figure}

\subsection{Độ phức tạp tính toán}
Trường hợp đơn giản nhất của bài toán -- gán $n$ tác vụ vào $m$ host với một
chiều tài nguyên duy nhất -- đã tương đương với bài toán \textit{Bin Packing
NP-hard}~\cite{garey1979computers}. Khi mở rộng sang nhiều chiều tài nguyên,
nhiều mục tiêu, và bản chất \textit{online} (tác vụ đến theo thời gian thực),
không thể trông đợi vào lời giải tối ưu chính xác trong thời gian đa thức. Đây
là động lực chính để chuyển sang các phương pháp dựa trên \textbf{học tăng
cường đa mục tiêu} (\textit{Multi-Objective Reinforcement Learning} -- MORL),
nơi agent học được chính sách (\textit{policy}) gần tối ưu thông qua tương tác
với môi trường mô phỏng~\cite{hayes2022practical, mao2016resource}.

\section{Mục tiêu đề tài}
\label{sec:muc-tieu}

\subsection{Mục tiêu tổng thể}
Đề tài \textbf{ELDAS} (\textit{Energy- and Latency-aware Datacenter AI
Scheduler}) hướng tới xây dựng một \textbf{nền tảng mô phỏng và huấn luyện
thuật toán MORL} cho bài toán lập lịch tác vụ AI trong trung tâm dữ liệu, với
khả năng:
\begin{itemize}
  \item Tái tạo trung thực hành vi của một cụm máy chủ AI ở quy mô vừa;
  \item Nạp được dữ liệu workload thực thay vì sinh ngẫu nhiên;
  \item So sánh công bằng nhiều thuật toán lập lịch (heuristic và RL) trên
        cùng một bộ kịch bản;
  \item Sản sinh các chỉ số định lượng đa mục tiêu phục vụ vẽ Pareto front.
\end{itemize}

\subsection{Mục tiêu Giai đoạn 1 (phạm vi báo cáo này)}
Báo cáo này trình bày sáu sản phẩm kỹ thuật của Giai đoạn 1. Mỗi sản
phẩm là một \emph{deliverable cụ thể, có thể kiểm chứng độc lập}:

\begin{enumerate}
  \item \textbf{MT1 -- Engine mô phỏng datacenter có theo dõi GPU.}
        Dựng môi trường mô phỏng cụm máy chủ trên nền CloudSim Plus,
        bổ sung lớp theo dõi GPU (mà CloudSim Plus không hỗ trợ sẵn)
        và mô hình công suất tuyến tính cho cả CPU lẫn GPU.
  \item \textbf{MT2 -- Tích hợp workload thực.} Parser đọc trực tiếp
        Alibaba PAI Cluster Trace v2023, kèm bộ lọc tạo ba kịch bản
        tải \emph{Low / High / Burst} để phục vụ so sánh.
  \item \textbf{MT3 -- Cầu nối Java--Python.} Dùng Py4J để Python
        điều khiển vòng đời mô phỏng (reset / step / pause / resume),
        đáp ứng đặc thù tương tác từng bước của RL.
  \item \textbf{MT4 -- Môi trường RL theo chuẩn MO-Gymnasium.}
        Phần thưởng giữ dạng \textbf{vector hai chiều}
        $(R^{\mathrm{energy}},\, R^{\mathrm{SLA}})$ -- không cộng dồn
        về một số; ràng buộc capacity được ép cứng qua action masking.
  \item \textbf{MT5 -- Hai scheduler baseline.}
        (i)~\emph{Random allocation} và (ii)~\emph{Kubernetes default}
        (Filter $\rightarrow$ LeastRequestedPriority). Hai baseline
        này tạo \emph{đường tham chiếu} cho agent MORL ở Phase 2 phải
        vượt qua.
  \item \textbf{MT6 -- Đóng gói tái sản xuất.} Toàn bộ hệ thống được
        gói trong Docker Compose; chạy cùng seed phải cho kết quả
        bit-for-bit trùng nhau, kể cả trên máy khác.
\end{enumerate}

\subsection{Mục tiêu Giai đoạn 2 (kế hoạch tiếp nối)}
Giai đoạn 2 dựa trên nền tảng MT1--MT6 để:
\begin{enumerate}
  \item Huấn luyện agent PPO + MLP, biến hai mục tiêu thành một bằng
        cách cộng có trọng số: $R = w \cdot R^{\mathrm{energy}} +
        (1 - w) \cdot R^{\mathrm{SLA}}$;
  \item Quét trọng số $w$ qua các giá trị
        $\{0{,}1,\, 0{,}2,\, \dots,\, 0{,}9\}$ để thu được nhiều điểm
        trên Pareto front;
  \item So sánh kết quả MORL với hai baseline trên cả ba kịch bản tải,
        đo trên ba chỉ số: năng lượng, makespan, tỉ lệ vi phạm SLA;
  \item Tinh chỉnh siêu tham số qua Weights \& Biases; thử Graph
        Neural Network thay MLP nếu thời gian cho phép.
\end{enumerate}

\section{Phạm vi và giới hạn}
\label{sec:pham-vi}

Để đảm bảo tính khả thi trong khuôn khổ một đồ án chuyên ngành, đề tài giới
hạn phạm vi như sau:

\paragraph{Về quy mô.}
Cụm mô phỏng cấu hình $5$ host (mỗi host gồm $32$ vCPU, $128$ GB RAM, $4$ GPU
NVIDIA A100), tải đầu vào trích từ Alibaba trace với kích thước
$1\,000$--$5\,000$ tác vụ. Quy mô này đủ để bộc lộ hiện tượng phân mảnh tài
nguyên và xung đột mục tiêu, đồng thời giữ thời gian huấn luyện một episode
dưới mức chấp nhận được trên phần cứng phòng thí nghiệm.

\paragraph{Về môi trường.}
Toàn bộ thực nghiệm được thực hiện trong môi trường \textbf{mô phỏng} bằng
CloudSim Plus, không triển khai trên cụm Kubernetes thực tế. Lý do: chu kỳ
huấn luyện RL đòi hỏi hàng triệu bước tương tác, không khả thi trên hạ tầng
thật.

\paragraph{Về phạm vi không bao gồm.}
Đề tài \textbf{không xét}: (i)~di chuyển máy ảo trực tiếp (\textit{live VM
migration}); (ii)~lập lịch xuyên trung tâm dữ liệu (\textit{multi-datacenter
scheduling}); (iii)~mô hình mạng chi tiết (\textit{network topology, congestion
modeling}); (iv)~lập lịch tác vụ phân tán có giao tiếp tập thể
(\textit{collective communication trong distributed training}).

\paragraph{Về tập hợp ràng buộc enforce trong Phase 1.}
Phát biểu bài toán ở §\ref{sec:bai-toan} liệt kê năm chiều
ràng buộc tài nguyên (CPU, RAM, GPU, băng thông, nhiệt độ). Phase 1
\textbf{chỉ thực thi cứng} ba chiều CPU/RAM/GPU
(§\ref{ssec:impl-baselines}); hai chiều còn lại được \emph{hoãn có chủ
đích} sang Phase 2 với lý do cụ thể như sau:

\begin{itemize}
  \item \textbf{Băng thông mạng.} Trace Alibaba PAI v2023 không cung
        cấp trường \texttt{bandwidth\_demand} ở mức task; bổ sung
        ràng buộc băng thông tại Phase 1 buộc phải sinh giả định tổng
        hợp (\textit{synthetic}, ví dụ tỉ lệ với \texttt{cpu\_milli}),
        làm giảm tính tái sản xuất khoa học của báo cáo và tạo ra
        một biến số không kiểm soát trong khi mục tiêu Phase 1 là
        \emph{kiểm chứng kiến trúc bốn lớp}. Hơn nữa, mỗi chiều ràng
        buộc thêm vào tăng kích thước observation thêm $H$ (một
        feature/host); với $H = 10$ điều này làm chậm hội tụ PPO
        không cần thiết khi chính thiết kế MORL còn đang được hiệu
        chỉnh. Phase 2 sẽ kích hoạt \texttt{host.getBwProvisioner()}
        sẵn có của CloudSim Plus (chỉ cần $\sim 20$ dòng code thêm
        vào \texttt{VmAllocationPolicy.find()}) sau khi đã có dữ liệu
        bandwidth thực hoặc một mô hình synthetic được biện luận chặt
        chẽ.
  \item \textbf{Nhiệt độ host.} Mô hình nhiệt yêu cầu (a)~mạng RC nhiệt
        ($T_j(t) = T_{\mathrm{amb}} + \theta_{\mathrm{JA}} \cdot P_j(t)$
        ở dạng đơn giản nhất), (b)~đường cong fan và ngưỡng throttle,
        và (c)~tương tác hai chiều với DVFS -- bản thân ba thành phần
        này là một nghiên cứu phụ. Quan trọng hơn, mô hình năng lượng
        Phase 1 \emph{tuyến tính} (§\ref{sec:nang-luong}) chưa có DVFS,
        nên việc thêm thermal lúc này sẽ là một \emph{lớp giả định trên
        lớp giả định} không có giá trị khoa học. Phase 2 dự kiến đưa
        DVFS + thermal vào cùng nhau như một gói nhất quán: trước hết
        thay \eqref{eq:power-model} bằng mô hình bậc hai có DVFS, sau
        đó thêm $T_j(t)$ làm feature observation và một
        \texttt{HostFitsThermal} filter dạng K8s predicate (cấm gán
        task lên host có $T_j > T_{\max}$).
\end{itemize}

Quyết định này tuân theo nguyên tắc \emph{minimum viable scope}: Phase 1
ưu tiên các ràng buộc \textbf{có dữ liệu thực} (CPU/RAM/GPU đều có trong
trace) để mọi con số trong báo cáo đều có thể truy xuất được nguồn;
các ràng buộc \textbf{cần dữ liệu giả định} được hoãn đến khi mô hình
nền (power, DVFS) đủ chính xác để không bị nhiễu giả định chồng lên
giả định.

\paragraph{Về mô hình năng lượng.}
Sử dụng mô hình tuyến tính \eqref{eq:power-model} cho cả CPU và GPU. Mô hình
phi tuyến (\textit{cubic, logarithmic}) và DVFS (\textit{Dynamic Voltage and
Frequency Scaling}) được ghi nhận như hướng mở rộng cho các giai đoạn sau.

\paragraph{Về thuật toán.}
Phase 1 chỉ hiện thực hai baseline (Random, K8s Default). Phase 2 tập trung
vào PPO + MLP với linear scalarization. Các phương pháp MORL phức tạp hơn như
Envelope Q-Learning, Pareto-DQN, hoặc các đầu Graph Neural Network được liệt
kê trong phần \textit{stretch goals}.

\section{Cấu trúc báo cáo}
\label{sec:cau-truc}

Báo cáo gồm bảy chương được tổ chức như sau:

\paragraph{Chương~\ref{chap:gioi-thieu} -- Giới thiệu.}
Trình bày bối cảnh, phát biểu hình thức bài toán lập lịch đa mục tiêu, mục
tiêu cụ thể, phạm vi và giới hạn của đề tài.

\paragraph{Chương 2 -- Cơ sở lý thuyết và nghiên cứu liên quan.}
Tổng quan ba mảng: (i)~thuật toán lập lịch truyền thống và Kubernetes
scheduler; (ii)~mô hình tiêu thụ năng lượng datacenter; (iii)~học tăng cường
đơn/đa mục tiêu với khảo sát ít nhất ba công trình tiêu biểu.

\paragraph{Chương 3 -- Kiến trúc hệ thống.}
Mô tả kiến trúc bốn lớp (Docker Compose, CloudSim Plus, Py4J Gateway,
MO-Gymnasium), các luồng dữ liệu chính và các quyết định thiết kế then chốt.

\paragraph{Chương 4 -- Hiện thực.}
Trình bày chi tiết hiện thực Phase 1: \texttt{SimulationConfig},
\texttt{DatacenterFactory}, \texttt{AlibabaTraceReader},
\texttt{SimulationManager} (lifecycle pause/resume/reset), Py4J
\texttt{GatewayEntryPoint}, hai \texttt{VmAllocationPolicy} baseline, và phía
Python: \texttt{state\_builder}, \texttt{reward}, \texttt{environment},
\texttt{tracker}.

\paragraph{Chương 5 -- Kết quả Phase 1.}
Báo cáo kết quả \textit{smoke test} vòng RL hoàn chỉnh, kiểm chứng tính đúng
của bridge Java--Python, đo các chỉ số ban đầu (năng lượng, SLA) của hai
baseline trên ba kịch bản tải.

\paragraph{Chương 6 -- Kế hoạch Phase 2.}
Lộ trình huấn luyện PPO, weight sweep, hyperparameter tuning và đánh giá
Pareto front; đi kèm sơ đồ Gantt và phân tích rủi ro.

\paragraph{Chương 7 -- Kết luận.}
Tổng kết đóng góp của Phase 1, bài học rút ra, và hướng phát triển.

Phần \textit{Phụ lục} cung cấp toàn bộ tệp cấu hình
(\texttt{docker-compose.yml}, \texttt{.env}), các đoạn mã đại diện và hướng
dẫn chạy lại thực nghiệm.
